\section{Introduction}

How well can a machine anticipate the moves of a professional investor?
Cohen, Lu, and Nguyen~\cite{cohen2026mimicking} find that LSTM models trained
on the prior positions of mutual fund managers can predict the direction of
their trades for the next quarter 71\% of the time, compared with 52\% for
each fund's most common action. Moreover, the managers and positions that are
not well predicted perform better. The next logical step is to test this with
hedge funds. The hedge funds that report to the SEC had \$5.4 trillion in net
assets (\$12.6 trillion gross) as of early 2025~\cite{sec2025pfstats}. They
are less constrained in their trading than mutual funds and are presumed to
trade on private information~\cite{agarwal2013hide, aragon2013confidential}.
Their only public record of stock positions is the quarterly 13F filing, and
many investors try to follow their trades. Time-series foundation
models~\cite{ansari2024chronos, ansari2025chronos2, das2024timesfm} have the
potential to make predictions without any training on the prediction task,
but their effectiveness in financial applications has been called into
question~\cite{rahimikia2025tsfm, noguer2026tsfm}.

We aim to answer the following questions:
\begin{itemize}
  \item \textbf{RQ1:} How predictable are hedge funds' quarterly trades (buy,
  hold or sell) using information that was public before the trades were
  disclosed? How does this compare with other managers and individual mutual
  funds using the same measure?
  \item \textbf{RQ2:} Can time-series foundation models or news improve the
  predictions?
  \item \textbf{RQ3:} Do unpredictable trades perform better?
\end{itemize}

Our results are mostly negative, but we designed the study to make the
negative results meaningful. For every null result, we report how large an
effect we can rule out. We address the most obvious loopholes: leakage,
survivorship bias in the selection of funds, weak or poorly chosen models,
class imbalance, and mismatched comparisons with other studies.

\paragraph{Contributions.}
\begin{enumerate}
  \item To the best of our knowledge, we are the first to apply the
  trade-direction design of~\cite{cohen2026mimicking} to hedge funds (prior
  work on predicting 13F holdings has instead predicted portfolio
  weights~\cite{izadifar2026holdings}). Using only free public SEC data, we run
  the same pipeline on five sets of funds: 49 well-known hedge funds, 99 hedge
  funds selected by rule from Form ADV filings made before the start of our
  sample period, 57 active long-only managers, and 300 active mutual funds from
  Form N-PORT, both individually and added up to their fund families.
  \item We find that hedge funds' trades are difficult to predict but not
  unusual: long-only managers and individual mutual funds are similarly
  difficult to predict in the same pipeline, and aggregation does not make
  them easier to predict. The high predictability found for mutual funds over
  1990--2023 does not hold for our period and data. With an oracle, we show
  that even when a fund's overall direction is known, its stock-level
  decisions remain largely unpredictable.
  \item We test three zero-shot foundation models (Chronos-2, TimesFM, and a
  Chronos model pre-trained only on financial data from before our test
  period), as well as news sentiment over several windows, against
  domain-trained models.
  \item We show with a positive control that our pipeline recovers
  predictability when it is planted, and we test, with multiple-testing
  correction and explicit power, whether unpredicted trades are more
  profitable before disclosure. We design the pipeline to avoid leakage and
  audit it by scrambling the inputs. Finally, we highlight a classic
  look-ahead bias that led us to overestimate our own early results:
  conditioning returns on a stock's future membership in an index.
\end{enumerate}
